% 来源：固定来源 PDF（2004福建文.pdf）第 2 页（印刷页 167）Q21 抛物线法线与割线中点几何示意图
\begin{tikzpicture}[scale=0.95, >=stealth, line cap=round, line join=round]
  % 坐标轴
  \draw[->, thin] (-3.2, 0) -- (3.2, 0) node[right] {$x$};
  \draw[->, thin] (0, -0.6) -- (0, 3.8) node[above] {$y$};
  \node[below left] at (0, 0) {$O$};

  % 抛物线 y = 0.5 x^2
  \draw[thick, domain=-2.6:2.6, samples=60] plot (\x, {0.5*\x*\x});

  % 选取代表性割线点 P 与 Q
  \coordinate (P) at (1.8, 1.62);
  \coordinate (Q) at (-1.4, 0.98);
  \coordinate (M) at ($(P)!0.5!(Q)$);

  % 割线 l
  \draw[thick] ($(Q)!-0.3!(P)$) -- ($(P)!-0.4!(Q)$) node[below right] {$l$};

  % 标注点
  \fill (P) circle (1.5pt) node[right] {$P$};
  \fill (Q) circle (1.5pt) node[above left] {$Q$};
  \fill (M) circle (1.8pt) node[above right] {$M$};
\end{tikzpicture}
